\documentclass[11pt,a4paper]{article}

\usepackage[utf8]{inputenc}
\usepackage[T1]{fontenc}
\usepackage[french]{babel}
\usepackage{lmodern}
\usepackage{geometry}
\usepackage{graphicx}
\usepackage{float}
\usepackage{amsmath}
\usepackage{amssymb}
\usepackage{array}
\usepackage{tabularx}
\usepackage{xcolor}
\usepackage{hyperref}
\usepackage{enumitem}
\usepackage{tcolorbox}

\geometry{margin=2.5cm}
\graphicspath{{../}{./}}

\hypersetup{
    colorlinks=true,
    linkcolor=blue!45!black,
    urlcolor=blue!45!black
}

\setlength{\parindent}{0pt}
\setlength{\parskip}{0.75em}
\setlist[itemize]{itemsep=0.25em, topsep=0.25em}

\definecolor{SoftGray}{RGB}{245,245,245}
\definecolor{DocRed}{RGB}{170,35,35}
\definecolor{DarkBlue}{RGB}{20,55,105}

\tcbset{
    colback=SoftGray,
    colframe=gray!45,
    arc=1mm,
    boxrule=0.4pt,
    left=2mm,
    right=2mm,
    top=1mm,
    bottom=1mm
}

\newcommand{\code}[1]{\texttt{#1}}
\newcommand{\COtwo}{CO$_2$}
\newcommand{\fluxunit}{\mathrm{W\,m^{-2}}}
\newcommand{\motcle}[1]{\textcolor{DarkBlue}{\textbf{#1}}}
\newcommand{\bloc}[1]{\vspace{0.9em}\noindent\textcolor{DarkBlue}{\underline{\textbf{#1}}}\par\vspace{0.35em}}
\newcommand{\titreprincipal}[1]{%
    \begin{tcolorbox}[
        colback=white,
        colframe=DocRed,
        boxrule=1pt,
        arc=1mm,
        left=3mm,
        right=3mm,
        top=2mm,
        bottom=2mm
    ]
    {\LARGE\bfseries\textcolor{DocRed}{#1}}
    \end{tcolorbox}
}
\newcommand{\partietitre}[1]{%
    \vspace{1.1em}
    {\Large\bfseries\textcolor{DocRed}{#1}}\par
    \vspace{0.35em}
}

\begin{document}

\begin{titlepage}
    \centering
    \vspace*{2cm}

    {\Huge\bfseries\textcolor{DocRed}{Document de Synthèse}\par}
    \vspace{0.8cm}
    {\LARGE\bfseries 2026-2027\par}
    \vspace{0.4cm}
    {\Large\par}

    \vfill

\end{titlepage}

\tableofcontents
\newpage

\titreprincipal{I -- Maintenance}
\addcontentsline{toc}{section}{I -- Maintenance}

\partietitre{1/Consigne et objectifs}
\addcontentsline{toc}{subsection}{Consigne et objectifs}

La partie maintenance vise à remettre en état le modèle hérité des années précédentes. L'objectif n'est pas de tout réécrire, mais de conserver un moteur utilisable, de clarifier les fichiers actifs et de séparer ce qui fonctionne de ce qui sert seulement de source pour la suite.

Le moteur stable par défaut reprend le modèle 4 des Carcajous. Les autres apports sont soit branchés explicitement, soit gardés comme références dans les documents sources.

\partietitre{2/ Organisation du dépôt}
\addcontentsline{toc}{subsection}{Organisation du dépôt}

La maintenance est regroupée dans le dossier \code{modele0\_maintenance}. Sa structure principale est la suivante :

\begin{itemize}
    \item \code{codes\_python/} : moteur, modules physiques et scripts de visualisation ;
    \item \code{ressources/} : données utilisées par le moteur ;
    \item \code{outils\_generation\_donnees/} : scripts de régénération des données ;
    \item \code{documents\_sources/} : PDF de synthèse récupérés des groupes précédents ;
    \item \code{PROVENANCE.md} et \code{THEORIE.md} : récapitualif de ce qui à été recupéré de quel groupe et rappels théoriques.
\end{itemize}

Le lancement de base se fait depuis \code{modele0\_maintenance/} avec :

\begin{verbatim}
python3 codes_python/modele_courbe.py --lat 48.5 --lon 2.3 --days 730
\end{verbatim}

\partietitre{3/Convection}
\addcontentsline{toc}{subsection}{Convection}

Deux formes de convection sont conservées dans le modèle maintenu :

\begin{itemize}
    \item la \motcle{convection forcée} provenant des Chevreaux ;
    \item la \motcle{convection naturelle} provenant des Ornithorynquietant.
\end{itemize}

Par défaut, les deux sont activées avec un vent constant de $2.5\ \mathrm{m\,s^{-1}}$. Elles peuvent être supprimées avec :

\begin{verbatim}
python3 codes_python/modele_courbe.py --sans-convection --no-plot
\end{verbatim}

\textcolor{green!50!black}{A faire: bien comprendre les 2 type de convection, quelle est leur distinction ? est qu'il y en à une des deux qui est négligeable ? moins pertinante ? Est ce que modéliser les 2 convection est juste physiquement ?}

\partietitre{4/Génération des données}
\addcontentsline{toc}{subsection}{Génération des données}

La génération des données passe par le script :

\begin{verbatim}
python3 outils_generation_donnees/generer_donnees.py
\end{verbatim}

Il permet de vérifier l'état des ressources, de régénérer les grilles rapides, les grilles annuelles, les CSV \code{12\_mois} et certaines données d'albédo. Les grilles rapides produites par défaut couvrent $7$ jours avec les deux convections actives.

Les données qui ne sont pas générées par code sont definies localement et restent à fournir ou remplacer manuellement.

\partietitre{5/Modules conservés}
\addcontentsline{toc}{subsection}{Modules conservés}

Certains modules sont gardés sans être totalement intégrés au moteur principal. C'est le cas de \code{codes\_python/physique/diffusion.py}, conservé pour une potentielle reprise future, mais pas active pour l'instant car mal comprise.

\partietitre{6/ Visualisation}
\addcontentsline{toc}{subsection}{Visualisation}

Les scripts de visualisation lisent les grilles déjà générées. Ils ne recalculent pas les données.

Exemples de commandes :

\begin{verbatim}
python3 codes_python/visualisation/modele_planisphere_haute_res.py --grille rapide
python3 codes_python/visualisation/modele_sphere_haute_res.py --grille rapide
python3 codes_python/visualisation/affichage_3D_rapide.py --month janvier --hour 12
\end{verbatim}

Les options de grilles disponibles sont \code{auto}, \code{rapide}, \code{1an} et \code{stabilisee}.

\partietitre{6/récapitulatif}
\addcontentsline{toc}{subsection}{Tableau récapitulatif}

\renewcommand{\arraystretch}{1.25}
\begin{table}[H]
\centering
\begin{tabularx}{\textwidth}{>{\bfseries}p{3.6cm} X}

\end{tabularx}
\caption{Résumé de la maintenance mettre le tableau de Pierre Louis}
\end{table}
\renewcommand{\arraystretch}{1}

\textcolor{green!50!black}{A faire: Inserer le tableau de Pierre Louis}

\newpage

\titreprincipal{II -- Modélisation de l'atmosphère}
\addcontentsline{toc}{section}{II -- Modélisation de l'atmosphère}

\partietitre{1/ Objectif et consigne}
\addcontentsline{toc}{subsection}{1/ Objectif et consigne}

L'objectif général est de modéliser l'effet du \COtwo{} présent dans l'atmosphère sur la puissance reçus et émis par la Terre.

À terme, ces flux pourront servir à calculer l'évolution de la température terrestre. Pour l'instant, on cherche surtout à comprendre comment le \COtwo{} modifie la puissance reçue et émise par la Terre.

Ce travail est mené en parallèle de la maintenance des modèles des années précédentes. L'objectif est d'intégrer progressivement cette modélisation de l'atmosphère au reste du projet CREPES.



\partietitre{2/ Modèle générale}
\addcontentsline{toc}{subsection}{2/ Modèle générale}

\bloc{Principe de découpage de l'atmosphère}

L'atmosphère est découpée suivant l'altitude en plusieurs couches. On suit ensuite les échanges radiatifs entre ces couches et la surface terrestre.

Pour simplifier, on considère que chaque couche possède :

\begin{itemize}
    \item une température uniforme, unique, indépendante du temps (dans les premiers modèles);
    \item une concentration en \COtwo{} uniforme et unique ;
    \item des propriétés radiatives moyennes à l'échelle de la couche.
\end{itemize}

On distingue les \motcle{couches} A, B, C etc, qui contiennent la matière atmosphérique, et les \motcle{interfaces}, où l'on suit les echanges de flux.

\begin{itemize}
    \item k=0 : surface, interface entre la Terre et la couche 1
    \item k=1 : interface entre la couche 1 et la couche 2
    \item ...
    \item k=N : interface entre la couche N et l'espace
\end{itemize}

\begin{figure}[H]
    \centering
    \includegraphics[width=0.72\textwidth]{Image_Docu.png}
    \caption{Schéma explicatif du principe de modélisation atmosphérique. On définit à la fois des couches atmosphériques et des interfaces où sont suivis les flux radiatifs.}
    \label{fig:schema-atmosphere}
\end{figure}

\textcolor{green!50!black}{A faire : Insérer le schéma réalisé par Melvin et Maud après avoir fais quelque modifs}

\bloc{Rayonnement visible et rayonnement infrarouge}

On distingue deux types de rayonnement :

\begin{itemize}
    \item le rayonnement solaire visible : emis par le Soleil ;
    \item le rayonnement infrarouge : emis par la Terre et les couches atmosphériques (car ce sont des corps noir cad des corps de température non nulle) rayonnement émis de manière isotrope(cad dans toutes les directions).
\end{itemize}

Les couches atmosphériques sont supposées \motcle{transparentes au visible}. Le rayonnement solaire traverse donc l'atmosphère. Dans l'infrarouge, en revanche, les couches peuvent absorber, transmettre et réémettre une partie du flux.

\renewcommand{\arraystretch}{1.35}
\begin{table}[H]
\centering
\begin{tabularx}{\textwidth}{>{\bfseries}p{3.2cm} X X}
\hline
 & Rayonnement lumineux visible & Rayonnement infrarouge \\
\hline
Couche atmosphérique
& Les couches sont transparentes au rayonnement visible : elles le laissent passer.
& Une couche atmosphérique, de température non nulle, émet un rayonnement infrarouge vers le haut et vers le bas. Elle peut aussi absorber une partie du rayonnement infrarouge reçu. \\
\hline
Terre
& La Terre est opaque : elle absorbe une fraction du rayonnement visible et réfléchit le reste selon son albédo.
& La Terre, de température non nulle, émet un rayonnement infrarouge. \\
\hline
\end{tabularx}
\caption{Comportement de la Terre et des couches atmosphériques vis-à-vis des rayonnements visible et infrarouge.}
\end{table}
\renewcommand{\arraystretch}{1}

\textcolor{green!50!black}{A faire : revoir le tableau récapitualitif (un peut redodant avec les explication du dessus), faire en sorte qu'il soit plus lisible (supprimer les longues phrase), ajouter des couleur pour rayonnement IR / visible en suivant le schéma du dessus}

\bloc{Flux solaire moyen absorbé par la surface}

La surface terrestre reçoit un flux solaire moyen. Une partie est réfléchie vers l'espace à cause de l'albédo, et le reste est absorbé.

On note :

\begin{itemize}
    \item $S_0$ le flux solaire total reçu par la Terre, en $\fluxunit$ ;
    \item $A$ l'albédo moyen ;
    \item $F_{\mathrm{SW}}$ le flux solaire moyen absorbé.
\end{itemize}

Le flux absorbé s'écrit :

\[
F_{\mathrm{SW}} = \frac{S_0(1-A)}{4}
\]

Le facteur $4$ vient du passage d'un disque éclairé à une moyenne sur toute la sphère terrestre : la Terre intercepte sur $\pi R^2$, puis on répartit sur $4\pi R^2$.

\bloc{Émission infrarouge totale de la surface}

La surface terrestre est traitée comme une surface opaque de température $T_s$. Comme tout corps de température non nulle, elle émet un rayonnement thermique.

Pour le flux total émis sur tout le spectre infrarouge, on utilise la loi de Stefan-Boltzmann :

\[
F_s = \sigma T_s^4
\]

où $\sigma$ est la constante de Stefan-Boltzmann :

\[
\sigma = 5.67\times10^{-8}\ \mathrm{W\,m^{-2}\,K^{-4}}
\]

Cette loi donne un \motcle{flux total} émis. Pour connaître la part de ce flux dans une bande absorbante du \COtwo{}, il faut utiliser la loi de Planck.

\bloc{Répartition spectrale : loi de Planck}

La loi de Planck donne la luminance spectrale d'un corps noir à la température $T$ :

\[
B_\lambda(T)
=
\frac{2hc^2}{\lambda^5}
\frac{1}{\exp\left(\frac{hc}{\lambda k_B T}\right)-1}
\]

avec :

\begin{itemize}
    \item $h$ la constante de Planck ;
    \item $c$ la vitesse de la lumière ;
    \item $k_B$ la constante de Boltzmann ;
    \item $\lambda$ la longueur d'onde ;
    \item $T$ la température du corps émetteur.
\end{itemize}

Pour obtenir le flux émis dans une bande spectrale $b=[\lambda_1,\lambda_2]$, on intègre la loi de Planck sur cette bande :

\[
E_b(T)
=
\int_{\lambda_1}^{\lambda_2}
\pi B_\lambda(T)\,d\lambda
\]

\textcolor{green!50!black}{Remarque : comprendre d'ou viens le pi et l'expliquer}

\textcolor{green!50!black}{A finir : finir de mettre au propre les calcul physique généraux valide pour tout les modèles}

\bloc{Absorbance, transmission et émissivité}


Le \COtwo{} n'absorbe pas uniformément tout l'infrarouge. Il absorbe surtout dans certaines bandes spectrales. Pour une bande $b$, on associe une absorbance moyenne $A_b$.

Cette absorbance est traitée comme une épaisseur optique effective. La transmission de la couche dans la bande $b$ est alors :

\[
\mathcal{T}_b = \exp(-A_b)
\]

Cette relation correspond à une forme de loi de Beer-Lambert : plus l'absorbance est grande, plus la transmission est faible.

L'émissivité effective de la couche dans la même bande est :

\[
\varepsilon_b = 1-\mathcal{T}_b
\]

Ce choix traduit la loi de Kirchhoff : à l'équilibre thermique local, une couche qui absorbe bien dans une bande émet aussi efficacement dans cette même bande.

Il ne faut pas confondre \motcle{absorbance} et \motcle{émissivité}. L'absorbance peut être supérieure à $1$. L'émissivité, elle, reste entre $0$ et $1$.

\bloc{Propagation du flux montant}

On note $F^\uparrow_{k,b}$ le flux infrarouge montant à l'interface $k$ dans la bande $b$.

À la surface, le flux montant est l'émission de la surface dans cette bande :

\[
F^\uparrow_{0,b} = E_b(T_s)
\]

Lorsqu'il traverse une couche, une partie du flux venant d'en bas est transmise, et la couche ajoute sa propre émission vers le haut :

\[
F^\uparrow_{k+1,b}
=
\mathcal{T}_{k,b} F^\uparrow_{k,b}
+
\left(1-\mathcal{T}_{k,b}\right)E_b(T_k)
\]

Le premier terme correspond au flux transmis. Le second correspond à l'émission de la couche $k$, pondérée par son émissivité dans la bande étudiée.

\bloc{Propagation du flux descendant}

On note $F^\downarrow_{k,b}$ le flux infrarouge descendant à l'interface $k$ dans la bande $b$.

Au sommet de l'atmosphère, l'espace n'apporte pas de flux infrarouge descendant dans ce modèle :

\[
F^\downarrow_{N,b} = 0
\]

En descendant vers la surface, chaque couche transmet une partie du flux venant d'en haut et ajoute sa propre émission vers le bas :

\[
F^\downarrow_{k,b}
=
\mathcal{T}_{k,b} F^\downarrow_{k+1,b}
+
\left(1-\mathcal{T}_{k,b}\right)E_b(T_k)
\]

Cette écriture est symétrique de celle du flux montant, mais appliquée dans le sens opposé.

\bloc{Flux transparent}

Toutes les longueurs d'onde qui ne sont pas traitées comme absorbantes sont regroupées dans un flux transparent. Ce flux est supposé traverser directement l'atmosphère.

Si $\mathcal{B}$ est l'ensemble des bandes absorbantes retenues, alors :

\[
F_{\mathrm{transparent}}
=
\sigma T_s^4
-
\sum_{b\in\mathcal{B}} E_b(T_s)
\]

Le flux montant total à une interface $k$ s'écrit donc :

\[
F^\uparrow_{k,\mathrm{total}}
=
F_{\mathrm{transparent}}
+
\sum_{b\in\mathcal{B}}F^\uparrow_{k,b}
\]

Le flux descendant total s'écrit :

\[
F^\downarrow_{k,\mathrm{total}}
=
\sum_{b\in\mathcal{B}}F^\downarrow_{k,b}
\]

Il n'y a pas de terme descendant transparent, car le modèle ne considère pas de source infrarouge venant de l'espace.

\bloc{Bilan radiatif}

À une interface, le flux net infrarouge est la différence entre le flux montant et le flux descendant :

\[
F_{k,\mathrm{net}}
=
F^\uparrow_{k,\mathrm{total}}
-
F^\downarrow_{k,\mathrm{total}}
\]

Pour une couche atmosphérique, le bilan énergétique complet comparerait ce qui entre dans la couche et ce qui en sort :

\[
C_k\frac{dT_k}{dt}
=
\sum_b
\left[
\left(F^\uparrow_{k,b}-F^\uparrow_{k+1,b}\right)
+
\left(F^\downarrow_{k+1,b}-F^\downarrow_{k,b}\right)
\right]
\]

Dans les premiers modèles, les températures restent imposées. Ce bilan sert donc surtout à préparer la suite : lorsque les températures ne seront plus fixées, il permettra de calculer leur évolution.

\partietitre{3/ Modèle 1}
\addcontentsline{toc}{subsection}{3/ Modèle 1}

Le modèle 1 est une première version simple du modèle général. Il sert a vérifie que nos calcul sont juste en cherchant à retomber sur une valeur connue au préalable.

Le script correspondant est : $modele1/modele1.py$

\bloc{Simplifications spécifiques au modèle 1}

En plus des hypothèses du modèle général, le modèle 1 ajoute les simplifications suivantes :

\begin{itemize}
    \item découpage en  trois couches atmosphériques ;
    \item les trois couches ont la même température et la même concentration en \COtwo{};
    \item seules deux bandes infrarouges du \COtwo{} sont traitées ;
    \item les absorbances des bandes sont fixées;
\end{itemize}

\bloc{Valeurs numériques utilisées}


\[
\begin{gathered}
T_s = 288.15\ \mathrm{K},
\qquad
F_s = \sigma T_s^4 = 390.9185\ \fluxunit,\\
S_0 = 1360\ \fluxunit,
\qquad
A = 0.30,
\qquad
F_{\mathrm{SW}} = 238\ \fluxunit,\\
T_{\mathrm{atm}} = 253.15\ \mathrm{K},
\qquad
C_{\mathrm{CO2}} = 425.65\ \mathrm{ppm}.
\end{gathered}
\]

Les trois couches sont \code{couche\_1} ($0$--$5\ \mathrm{km}$), \code{couche\_2} ($5$--$10\ \mathrm{km}$) et \code{couche\_3} ($10$--$20\ \mathrm{km}$). Elles ont toutes la même température et la même concentration en \COtwo.

Les deux bandes retenues sont \code{CO2\_15um}, de $14.25$ à $15.75\ \mu\mathrm{m}$ avec $A_b=1.0$, et \code{CO2\_4\_3um}, de $4.20$ à $4.35\ \mu\mathrm{m}$ avec $A_b=3.25$.

\[
\begin{aligned}
A_b=1.0 &\Rightarrow \mathcal{T}_b=0.3679,\quad \varepsilon_b=0.6321,\\
A_b=3.25 &\Rightarrow \mathcal{T}_b=0.0388,\quad \varepsilon_b=0.9612.
\end{aligned}
\]

Pour la surface, les émissions par bande valent :

\[
E_{\mathrm{CO2,15\mu m}}(T_s)=27.4827\ \fluxunit,
\qquad
E_{\mathrm{CO2,4.3\mu m}}(T_s)=0.3331\ \fluxunit
\]

Le flux transparent est donc :

\[
F_{\mathrm{transparent}}
=
390.9185-(27.4827+0.3331)
=
363.1027\ \fluxunit
\]

\textcolor{green!50!black}{A faire : bien comprendre les valeur numerique utilisées, sourcer et peut etre revoir la mise en page}

\bloc{Résultat du modèle 1}

La commande d'exécution est : $python modele1/modele1.py$

La sortie actuelle est :

\begin{verbatim}
flux_infrarouge_sortant_sommet_atmosphere_W_m2 = 380.782258
flux_infrarouge_descendant_surface_W_m2 = 16.311241
\end{verbatim}

Le premier résultat est le \motcle{flux infrarouge sortant} au sommet de l'atmosphère. Le second est le \motcle{flux infrarouge descendant} reçu par la surface.

\bloc{Vérification limite}

On vérifie le comportement du modèle en annulant les absorbances. Si $A_b=0$, alors $\mathcal{T}_b=1$ et $\varepsilon_b=0$ : les couches deviennent transparentes et n'émettent plus dans les bandes étudiées.

Le résultat attendu est donc :

\[
F^\downarrow_{\mathrm{surface}} = 0,
\qquad
F^\uparrow_{\mathrm{sommet}} = \sigma T_s^4
\]

Le test donne :

\begin{verbatim}
flux_infrarouge_sortant_sommet_atmosphere = 390.918507769 W m^-2
flux_infrarouge_descendant_surface = 0 W m^-2
\end{verbatim}

Le test est cohérent : sans absorption, il n'y a pas de réémission atmosphérique.

\partietitre{4/ Modèle 2}
\addcontentsline{toc}{subsection}{4/ Modèle 2}

\textcolor{green!50!black}{A faire}

\end{document}
